\section{Trap Bandit Experiments}
\label{section:trap-bandits}

We consider a simple experiment illustrating how a robust infra-Bayesian learner can be useful even in a stateless stochastic bandit setting. The environment has two arms, indexed by \(i \in \{1,2\}\). At the start of each run, the arm reward probabilities \((p_1,p_2)\) are sampled uniformly from
\[
    \{(0.3,0.7), (0.7,0.3)\},
\]
so that one arm has significantly higher expected reward than the other. Independently, the world is sampled to be either safe or risky, with probability \(\alpha\) of being risky and probability \(1-\alpha\) of being safe.

In a safe world, each arm is Bernoulli with reward probability \(p_i\). In a risky world, the arm with the larger reward probability is also the trapped arm. Pulling this arm yields a catastrophic reward of \(-1000\) with probability \(p_{\mathrm{cat}}\), yields reward \(1\) with probability \(p_i\), and yields reward \(0\) otherwise. The lower-probability arm remains Bernoulli. Thus each run is generated as follows:
\[
\begin{array}{ll}
\text{sample} & (p_1,p_2) \sim \mathrm{Unif}\{(0.3,0.7),(0.7,0.3)\},\\
\text{sample} & \text{world type} \sim \mathrm{Bernoulli}(\alpha),\\[2mm]
\text{if safe:} & r_i \sim \mathrm{Bernoulli}(p_i),\\[1mm]
\text{if risky:} &
\begin{cases}
\text{the trap arm } \arg\max_i p_i \text{ yields } -1000 \text{ with probability } p_{\mathrm{cat}},\\
\text{the trap arm yields } 1 \text{ with probability } p_i,\\
\text{the trap arm yields } 0 \text{ otherwise},\\
\text{the other arm yields } r_i \sim \mathrm{Bernoulli}(p_i).
\end{cases}
\end{array}
\]

We compare classical Bayesian agents with an infra-Bayesian agent using the same joint hypothesis machinery. Bayesian agents represent uncertainty using \texttt{Infradistribution.mix(...)}. The infra-Bayesian agent instead uses Knightian uncertainty over the safe and risky world families via \texttt{Infradistribution.mixKU(...)}, while retaining ordinary Bayesian uncertainty over \((p_1,p_2)\) within each family.

The experiments vary the relationship between the Bayesian point prior on the risky-world probability and the true data-generating process. We set $\alpha$ of the data-generating process to be $0.99$ in our experiments. In the first condition, the Bayesian prior on \(P(\mathrm{risky})\) is correctly specified in expectation. We then consider misspecified point-prior conditions in which Bayesian agents assign too little probability to the risky world ($\alpha=0.01$). Across all conditions, the infra-Bayesian agent uses the same classical prior over \((p_1,p_2)\) as the Bayesian agents, but maintains Knightian uncertainty over whether the world is safe or risky.

For Bayesian agents, we evaluate two exploration strategies: greedy action selection and Thompson sampling. The infra-Bayesian agent uses greedy action selection with respect to its robust lower values, with uniform tie-breaking. Regret is measured relative to the best policy with full knowledge of the true world. We report cumulative expected regret percentiles and trapped-arm pull-rate percentiles. Results are shown in Figure \ref{fig:trapped-bandits}.

\begin{figure}[htb]
    \centering
    \includegraphics[width=\textwidth]{trapped_bandits.png}
    \caption{Comparing the performance of infra-Bayesian and classical Bayesian agents (with either greedy or Thompson Sampling exploration strategies) in the trap bandit setting. The first row shows results for a correctly specified Bayes prior condition; the second for a severely misspecified Bayes prior condition. The first column shows cumulative expected regret, and the second shows the average pull rate of the risky (trap) arm.}
    \label{fig:trapped-bandits}
\end{figure}

We observe that the greedy Bayesian and infra-Bayesian agents perform identically in the risky world when the Bayesian prior is properly specified. Notably, the infra-Bayesian is able to properly learn the classical dynamics over arms, which it must determine to be able to act conservatively, even when maintaining Knightian uncertainty at the topmost layer (as to whether it is in a risky world). It makes sense that the properly specified Bayesian agent behaves this way because it is favorable expected value to be conservative when the likelihood of being in a risky world is so high. However, if the prior is severely misspecified, and the agent believes they are in a mostly safe world where it would pay to exploit the high-reward arm, they suffer catastrophic regret. Thompson sampling is slightly helpful in fixing the poor prior the misspecified prior setting; it is subtly detrimental (as expected at short time horizons) in the correctly specified prior setting. 

These results demonstrate Infra-Bayesian agents matching or outperforming classical agents in risky worlds, which raises a natural question: at what cost? Table \ref{tab:trap-bandit-results} includes an alterative mostly-safe scenario, where the data generating process sets $\alpha$ to $0.01$; in other words, the world is generally safe. In this setting, the infra-Bayesian agent incurs substantially higher regret, reflecting the cost of maintaining Knightian uncertainty over the risky-world hypothesis. 

\begin{table}[t]
\centering
\small
\begin{tabular}{lll rcc}
\toprule
$\alpha$ (DGP) & $\alpha$ (Prior) & Agent & Cat. rate & p50, 95\% CI & p95, 95\% CI \\
\midrule
0.99 & n/a & infra\_bayesian & 0.045 & 9.60 [9.60, 9.60] & 125.76 [105.60, 183.36] \\
0.99 & 0.99 & bayes\_greedy & 0.045 & 9.60 [9.60, 9.60] & 125.76 [96.48, 183.36] \\
0.99 & 0.01 & bayes\_greedy & 0.645 & 624.00 [528.00, 748.80] & 960.00 [950.88, 960.00] \\
0.99 & 0.99 & bayes\_thompson & 0.080 & 38.40 [28.80, 38.40] & 115.20 [96.00, 144.00] \\
0.99 & 0.01 & bayes\_thompson & 0.645 & 580.80 [489.60, 720.00] & 950.40 [940.80, 950.40] \\
0.01 & n/a & infra\_bayesian & 0.000 & 39.60 [39.60, 39.60] & 40.00 [40.00, 40.00] \\
0.01 & 0.01 & bayes\_greedy & 0.015 & 0.40 [0.40, 0.80] & 5.60 [4.02, 9.64] \\
0.01 & 0.01 & bayes\_thompson & 0.015 & 1.60 [1.20, 1.60] & 7.60 [5.22, 11.60] \\
\bottomrule
\end{tabular}
\caption{Final cumulative expected-regret percentiles with bootstrap confidence intervals.}
\label{tab:trap-bandit-results}
\end{table}

